Le skeleton est un sport d'hiver qui se pratique individuellement sur une
planche où le skeletoneur se place à plat ventre, la tête devant. L'objectif est
de parcourir la piste le plus rapidement possible. On se propose d'étudier le
mouvement du centre d'inertie $G$ d'un skeletoneur, de masse $m$, sur une partie
d'une piste rectiligne inclinée d'un angle $\alpha$ par rapport à l'horizontal.
L'étude est réalisée dans un repère ($O,\ex,\ey$ ) liée à la Terre
supposé galiléen.

\begin{minipage}{0.65\linewidth}
    On choisit l'instant de passage de $G$ en $O$ comme origine de temps
    ($t_{0}=0$). On repère la position de $G$ à un instant $t$ par son abscisse
    $x_{G}$ dans ce repère (Figure~\ref{fig:skeletoneur}).
    
    Le skeletoneur arrive au point de départ $O$ avec la vitesse
    $\vec{v}_{0}=v_{0}\ex$, il passe ensuite par la position $A$ avec une
    vitesse $\vec{v}_{A}$.
    
    Le contact avec le plan incliné se fait avec frottement modélisé par une
    force $\vec{f}$ parallèle au plan, de sens opposée au sens de mouvement et
    d'intensité constante $f$.
    
    \begin{donnees}
        \begin{itemize}
            \item $g=\qty{9,8}{\a}$~; $\alpha=\ang{10}$~; $m=\qty{80}{kg}$~;
                  $v_{A}=\qty{30}{\v}$.
        \end{itemize}
    \end{donnees}
\end{minipage}
\hfill
\begin{minipage}{0.34\linewidth}
    \begin{figure}[H]
        \centering
        \includegraphics[width=\textwidth]{2025-05-17_014726-}
        \caption{}
        \label{fig:skeletoneur}
    \end{figure}
\end{minipage}

\begin{minipage}{0.65\linewidth}
\begin{enumerate}
    \item En appliquant la deuxième loi de Newton, montrer que l'accélération de
          $G$ s'écrit~: $a_{G}=g \sin \alpha-\frac{f}{m}$.
    \item La courbe de la figure~(\ref{fig:v_de_t}) donne l'évolution de la
          vitesse $v$ de $G$ au cours du temps.
          \begin{enumerate}
              \item Déterminer la valeur de $a_{G}$.
              \item Déduire, en justifiant, la nature du mouvement de $G$.
          \end{enumerate}   
    \item Calculer la valeur de $f$.
    \item Déterminer la valeur de l'instant $t_{A}$ de passage de $G$ par $A$.
    \item Calculer la distance $OA$.
\end{enumerate}
\end{minipage}
\hfill
\begin{minipage}{0.34\linewidth}
    \begin{figure}[H]
        \centering
        \includegraphics[width=\linewidth]{2025-05-17_015326-}
        \caption{}
        \label{fig:v_de_t}
    \end{figure}
\end{minipage}

